% !PRINTED 512
\weberpage{512}

From this we further conclude, by formula (1), that \gesperrt{associated functionals have the same absolute norm}.

According to this, an integral rational functional is always associated with its absolute value. If then $\alpha$ is any integral functional of the field $\overline{\Omega}$, then $N(\alpha)$ and $N_a(\alpha)$ are also associated. Now since in $N(\alpha)=\alpha\alpha'$ the factor $\alpha'$, being a product of integral functionals, is itself integral, $N(\alpha)$, and consequently also the natural number $N_a(\alpha)$, is divisible by $\alpha$. We therefore formulate the further theorem:

13. \gesperrt{There exist natural integers (in infinite number), among them the absolute norm of $\alpha$, which are divisible by an arbitrary integral functional $\alpha$. If $a$ is the least among the natural numbers divisible by $\alpha$, then every integral rational number divisible by $\alpha$ is divisible by $a$.}

For if $m$ is an integral rational number divisible by $\alpha$, then the residue of the division of $m$ by $a$ is also an integral rational number divisible by $\alpha$. Since this is smaller than $a$, it must be $=0$.

\S.\,139.

\gesperrt{Greatest common divisor.}

Let $\alpha, \beta, \ldots$ denote integral functionals in $\Omega$ different from zero, in arbitrary but finite number, and $x, y, \ldots$ variables which do not occur in $\alpha, \beta, \ldots$. Then

\weq{1}{\delta = \alpha x + \beta y + \cdots}

is likewise an integral functional in $\Omega$. The norm of $\delta$ is an integral homogeneous function of degree $n$ of the variables $x, y, \ldots$, whose coefficients are integral rational functionals. If we denote the absolute norm of $\delta$ by $D$, then

\weq{2}{N(\delta) = D\,E(x, y, \ldots) = D\,E,}

and in this $E(x, y, \ldots) = E$ is an integral rational function of the variables $x, y, \ldots$ and in general a fractional function of the variables occurring in $\alpha, \beta, \ldots$, but in any case a \gesperrt{rational unit} [\S.\,136, (2)].
